\section{Architecture：原生多模态与视觉代理}

\subsection{从 KiMi VL 到 K2.5}
KiMi 早期保持纯语言模型与 VL 模型分离，但到了 K2.5 开始就不再区分，讲师用了 ``原生多模态'' 来描述那种从 day 0 就把视觉 token、工具调用以及文本融合的架构。K2.5 不再是 ``先生成语言再塞一帧图''，而是让视觉、文本、工具在同一个 transformer 中共学，把 attention 视作一个 multi-agent 的通道。

\begin{importantbox}{原生多模态的演进}
\begin{itemize}
    \item 旧版本：K2、KiMi VL、Qwen3-VL 会把语言模型和视觉模型割裂，靠 late fusion 实现图文对应。
    \item K2.5：从第 0 天开始把图片、视频帧、命令行以及工具调用都拼在一起，early fusion 让每一层都有联合模态信号。
    \item 这意味着 ``视觉'' 不是外加能力，而是贯穿整个 transformer attention 的核心维度，prompt 里甚至会写明 ``视觉 token 先于语言 token 出现''。
\end{itemize}
\end{importantbox}

\subsection{Tokenization、Memory 与 Early Fusion}
K2.5 进一步把 tokenization 和 memory 设计成多模态感知：视觉帧被切成 patch tokens，并加上 temporal position，再与语言 token 混合进入 cache。讲师提到 early fusion 并非把所有 token 一起扔进 attention，而是先用 `modality router` 做 soft gate，确保视觉信息在每一层都得到重采样，而不会因为文本权重过大被掩盖。

\begin{knowledgebox}{Early fusion 的三层策略}
\begin{itemize}
    \item 用 modality routing 来平衡视觉/语言 token，在同一 layer 里维持 2:1 的 visual-text ratio。
    \item 把 temporal memory cache 设计成 `frame bundle + tool trace` 的结构，方便 later evaluation 回放一个 agent loop。
    \item 把 attention mask 设计成 ``vision first''，让视觉 token 能先冲进 GRM 评分 pipeline，再决定是否继续下一步。
\end{itemize}
\end{knowledgebox}

\subsection{Visual Agency Intelligence 的自组织}
K2.5 的核心定位是 ``Visual Agency Intelligence''：不仅理解视觉内容，更能启动 agent loop 去执行任务、调用工具甚至控制 GUI。讲师提到 `Activation the Vision-Agentic Capability`（在向导中那条脉络），说明模型在前向推理时要同时看到视觉、语言与工具的 token，然后一边感知、一边提出下一步的操作。

\begin{figure}[H]
\centering
\includegraphics[width=0.78\textwidth,trim=0 120 0 120,clip]{cover.jpg}
\caption{封面画面里同时出现讲师、Visual Agentic 架构图和 K2.5 timeline，强调视觉与 agent loop 的联合演讲。\protect\footnotemark}
\end{figure}
\footnotetext{视频画面时间区间：00:32:10--00:32:48，画面里讲师展示 agent swarm 与原生多模态的叠加。}

\subsection{Agent loop 与 GUI agents}
在 architecture 层面，visual agency 不只停留在 attention 里，而是通过 agent loop 接管实际行动：模型会用 VRL（Visual Reinforcement Loop）规划，先生成 intent、再调用工具、最后回到 Allcastrater。GUI agents 是这一层面的延伸，它们专门模拟界面交互，确保 architecture 能用标准化 prompt 控制按钮、滑块、表格甚至浏览器标签页。

\begin{knowledgebox}{GUI agents 的任务边界}
\begin{itemize}
    \item GUI agents 担任界面感知与操作：把 screenshot 解析成 mouse/key events。
    \item 工具 agent 负责 API 调用：在 API response 里提取 state、写入 tool trace。
    \item 文本 agent 则处理解释性的输出，保证 GRM 能从多模态输入中统一评分。
\end{itemize}
\end{knowledgebox}

\subsection{解释性与信号路由}
原生多模态让 attention layers 同时观测图像、文本、工具 token，方便我们追踪模型做出的每一步。这条路径也服务于 explainability，比如在训练日志里标出 ``vision tracker''、``tool calling track''，把 reward 信号、操作意图、GRM 得分串联起来，帮助我们在 model debug 时快速定位问题。讲师特别提到，用 GRM scores + tool trace 组成 timeline，可以在 10 秒内复现某个 agent 随机生成的 action。

\subsection{本章小结}
K2.5 的架构把以前的语言与视觉割裂彻底打通，early fusion 和 agent-centric design 让 ``Visual Agency'' 成为模型的默认工作方式，像 figure 里强调的那样：视觉、语言、工具三者同步起舞。
